\section{Longest Bitonic Subsequence}
\label{sec:dp-longest-bitonic-subsequence}

\problemstatement{Given an array $\textit{arr}$, find the length of the
longest \textbf{bitonic} subsequence: one that first strictly increases,
then strictly decreases (either the increasing or decreasing part may
be empty, but not both).}

\begin{patternbox}
\emph{``Increases then decreases''} is solved by running
Section~\ref{sec:dp-lis}'s ending-here LIS computation \textbf{twice} --
once forward (for the increasing run up to each index) and once
backward (for the decreasing run from each index onward) -- then
combining the two values at every index, since a bitonic sequence's
peak can be at any position.
\end{patternbox}

\subsection*{Deriving the combination formula}

Let $\textit{inc}[i]$ = length of the longest increasing subsequence
\emph{ending} at index $i$ (exactly Section~\ref{sec:dp-lis}'s table).
Let $\textit{dec}[i]$ = length of the longest decreasing subsequence
\emph{starting} at index $i$ -- computed identically, but scanning
\emph{backward} from the end and requiring $\textit{arr}[i] >
\textit{arr}[j]$ for $j>i$ (equivalently, LIS run on the array's mirror
image). If index $i$ is the peak of a bitonic sequence, the increasing
run contributes $\textit{inc}[i]$ elements and the decreasing run
contributes $\textit{dec}[i]$ elements, but $\textit{arr}[i]$ itself is
counted in both, so:
\[
  \textit{bitonic}[i] = \textit{inc}[i] + \textit{dec}[i] - 1.
\]
The answer is $\max_i \textit{bitonic}[i]$.

\subsection*{C++ implementation}

\begin{lstlisting}
#include <bits/stdc++.h>
using namespace std;

int longestBitonicSubsequence(vector<int>& arr) {
    int n = arr.size();
    vector<int> inc(n, 1), dec(n, 1);

    for (int i = 0; i < n; i++) {
        for (int j = 0; j < i; j++) {
            if (arr[j] < arr[i]) {
                inc[i] = max(inc[i], 1 + inc[j]);
            }
        }
    }

    for (int i = n - 1; i >= 0; i--) {
        for (int j = n - 1; j > i; j--) {
            if (arr[j] < arr[i]) {
                dec[i] = max(dec[i], 1 + dec[j]);
            }
        }
    }

    int best = 0;
    for (int i = 0; i < n; i++) {
        best = max(best, inc[i] + dec[i] - 1);
    }
    return best;
}

int main() {
    vector<int> arr = {1, 11, 2, 10, 4, 5, 2, 1};
    cout << longestBitonicSubsequence(arr) << endl; // 6
    return 0;
}
\end{lstlisting}

\subsection*{Dry run}

\dryrun{Take $\textit{arr} = [1,11,2,10,4,5,2,1]$.}

$\textit{inc} = [1,2,2,3,3,4,2,1]$ (the run $1,2,4,5$ ending at index
$5$ gives $\textit{inc}[5]=4$). $\textit{dec} = [1,5,2,4,3,3,2,1]$ (the
run $11,10,5,2,1$ starting at index $1$ gives $\textit{dec}[1]=5$).
$\textit{bitonic}[1] = \textit{inc}[1]+\textit{dec}[1]-1 = 2+5-1=6$,
the maximum -- realized by the sequence $1,11,10,5,2,1$, increasing
from $1$ to the peak $11$ (index $1$) and then decreasing.

\begin{complexitybox}
\textbf{Time Complexity.} $O(n^2)$ -- two nested-loop passes.
\textbf{Space Complexity.} $O(n)$ for \textit{inc} and \textit{dec}.
\end{complexitybox}

\begin{takeawaybox}
When a problem's structure splits naturally into a ``before the
turning point'' half and an ``after the turning point'' half, compute
each half's ending-here table independently (one forward, one
backward) and combine them at every candidate turning point -- a
pattern that generalizes well beyond bitonic sequences to any
problem with a single peak or pivot.
\end{takeawaybox}
